\section{The Order the Graph Decides}
\label{sec:ch12-the-order-the-graph-decides}
\index{extraction!order of|(}%

There is no interesting question about \emph{whether} the five remaining domains
could be extracted. Chapter~\ref{ch:the-first-cut-extracting-catalog} did one and
the rest are the same shape. The interesting question is the order, and the graph
answers it if you ask it properly.

\index{domain dependencies!read off the type files}%
Asking it properly means reading the type files rather than the folder names. A
dependency that matters here is not a \code{using} line; it is a field that
returns a type another domain owns, because that is the only kind of coupling
that survives a process boundary. Eight classes in \code{Mosaic.Api} carry an
\code{[ObjectType<T>]} attribute at tag \code{ch11}. Seven of them are domain
types and the whole answer is in those; the eighth is the empty partial that
marks \code{PageCursor} shareable.

\begin{minted}[fontsize=\footnotesize,breaklines]{text}
Pricing/Types/ProductPricingNode.cs      [ObjectType<Product>]
Inventory/Types/ProductInventoryNode.cs  [ObjectType<Product>]
Reviews/Types/ProductReviewsNode.cs      [ObjectType<Product>]
Reviews/Types/ReviewNode.cs              [ObjectType<Review>]
Accounts/Types/CustomerNode.cs           [ObjectType<Customer>]
Ordering/Types/OrderNode.cs              [ObjectType<Order>]
Ordering/Types/OrderLineNode.cs          [ObjectType<OrderLine>]
\end{minted}

Three of them hang a field off \code{Product}, and hanging a field off a type is
not a dependency on the domain that owns it:
chapter~\ref{ch:the-first-cut-extracting-catalog}'s extraction proved that by
moving Catalog out from under all three without any of them changing by a
character. A dependency is a field whose \emph{return type} belongs elsewhere,
and there are exactly three of those:

\begin{minted}[fontsize=\footnotesize,breaklines]{text}
Review.author       -> Customer     Reviews  depends on Accounts
Order.customer      -> Customer     Ordering depends on Accounts
OrderLine.product   -> Product      Ordering depends on Catalog
\end{minted}

So Pricing, Inventory and Accounts reference nobody, Reviews references
Accounts, and Ordering references both Accounts and Catalog.
Figure~\ref{fig:ch12-domain-dependencies} puts the two states side by side, and the
thing to notice about it is that the arrows do not move. The same three
references exist before and after. What changes is what an arrow costs.

\begin{figure}[htbp]
  \centering
  % Which Mosaic domain references which, before and after the extraction.
% Referenced from chapter 12. Bare tikzpicture; the figure environment, caption
% and label live at the call site.
%
% The point of the picture is that the arrows do not move. The same three
% references exist on both sides; what changes is that on the left they are
% method calls inside one process and on the right they are keys crossing a
% network. The extraction order falls straight out of the arrow directions:
% a box nothing points away from can go first.
\begin{tikzpicture}[
    font=\small,
    head/.style={font=\small\bfseries, align=center},
    dom/.style={draw, rounded corners=2pt, minimum width=19mm,
                minimum height=5.5mm, align=center, font=\scriptsize},
    svc/.style={draw, rounded corners=2pt, minimum width=19mm,
                minimum height=5.5mm, align=center, font=\scriptsize,
                fill=black!8},
    proc/.style={draw, dashed, rounded corners=2pt, gray},
    ref/.style={-{Stealth[length=2mm]}},
    note/.style={font=\scriptsize, gray, align=center, inner sep=1pt}
  ]

  % -- left: one process, five domains, plus Catalog next door ---------------
  \node[head] at (2.4,0.8) {at \code{ch11}: a method call};

  \draw[proc] (0.15,-3.45) rectangle (4.65,0.3);
  \node[note, anchor=north west] at (0.2,-3.5) {one process};

  \node[dom] (pri) at (1.4,-0.1)  {Pricing};
  \node[dom] (inv) at (3.4,-0.1)  {Inventory};
  \node[dom] (acc) at (1.4,-1.2)  {Accounts};
  \node[dom] (rev) at (3.4,-1.2)  {Reviews};
  \node[dom] (ord) at (2.4,-2.4)  {Ordering};
  \node[svc] (cat) at (2.4,-4.3)  {Catalog};

  \draw[ref] (rev) -- (acc);
  \draw[ref] (ord) -- (acc);
  \draw[ref] (ord) -- (cat);

  % -- right: six processes, the same three arrows ---------------------------
  \node[head] at (9.6,0.8) {at \code{ch12}: a key and a hop};

  \node[svc] (pri2) at (8.6,-0.1)  {Pricing};
  \node[svc] (inv2) at (10.6,-0.1) {Inventory};
  \node[svc] (acc2) at (8.6,-1.2)  {Accounts};
  \node[svc] (rev2) at (10.6,-1.2) {Reviews};
  \node[svc] (ord2) at (9.6,-2.4)  {Ordering};
  \node[svc] (cat2) at (9.6,-4.3)  {Catalog};

  \draw[ref] (rev2) -- (acc2);
  \draw[ref] (ord2) -- (acc2);
  \draw[ref] (ord2) -- (cat2);

  % -- the extraction order, read off the arrows -----------------------------
  \draw[gray] (7.2,-0.05) -- (6.95,-0.05) -- (6.95,-1.25) -- (7.2,-1.25);
  \node[note, anchor=east, align=right] at (6.85,-0.65)
    {reference nobody:\\go first};
  \draw[gray] (7.2,-1.6) -- (6.95,-1.6) -- (6.95,-2.45) -- (7.2,-2.45);
  \node[note, anchor=east, align=right] at (6.85,-2.05)
    {reference somebody:\\go last};

\end{tikzpicture}

  \caption{The same three references, twice. On the left they are method calls
  inside one process, resolved by a DataLoader against a table in the same
  database. On the right they are keys in wrappers, resolved by the router
  against another service. Nothing about the domain model changed; the
  extraction order is read straight off the arrow directions, because a box
  nothing points away from can leave without anybody noticing.}
  \label{fig:ch12-domain-dependencies}
\end{figure}

\index{extraction!leaves first}%
The order that falls out is Pricing, Inventory, Accounts, then Reviews, then
Ordering. A domain that references nobody can be lifted out on its own: nothing
in the service it leaves has to change, because nothing in the service it leaves
was calling it. A domain that references somebody has to go after the domain it
references, or it goes out holding a compile-time reference to a project that is
about to become a network hop.

That is the rule chapter~\ref{ch:hotchocolate-quickly} wrote down as a layout
convention and has been holding ever since, and this chapter is where it gets
paid. I want to be honest that the payment is smaller than the rule makes it
sound. Extracting the three leaves cost nothing because there was nothing to
untangle; extracting Reviews and Ordering cost exactly the three fields above,
and section~\ref{sec:ch12-what-the-boundary-costs} is where the bill for those
three arrives.

\subsection*{Six databases, and one server}

\index{database!one per service}%
Each service takes a database of its own on the PostgreSQL container all six
share. Chapter~\ref{ch:the-first-cut-extracting-catalog} did this once for
Catalog and called it a deployment detail. Doing it five more times is where a
detail turns into a policy, and the policy is that two services never share a
table.

Nothing in the compose file creates any of them. \code{EnsureCreatedAsync} does,
on each service's first start, which is why this extraction and
chapter~\ref{ch:the-first-cut-extracting-catalog}'s between them added six
services to \code{docker-compose.yml} and no storage at all:

\begin{minted}[fontsize=\footnotesize]{text}
catalog    5101   database catalog
pricing    5102   database pricing
inventory  5103   database inventory
accounts   5104   database accounts
reviews    5105   database reviews
ordering   5106   database ordering
\end{minted}

Six services on one server is a convenience I would not defend in production and
would not spend a chapter apologising for either. What is being defended is the
line above it: no service can read another's tables, so no shortcut across a
seam is one connection string away. The credentials are still the shared
development ones, and that would be a finding anywhere but here.

There is one number in that table worth pausing on, and it is the one that is
missing. Nothing listens on 5100.

\medskip
That is the plan. What it costs to carry out is a different question, and the
answer is mostly boring, which is the good outcome.

\index{extraction!order of|)}%
